\section{A threshold proof for fixed quadratic padding}
\label{c:rat-threshold}

The preceding overlap proof assumes that both labels fail and then
compares their heights.  For fixed quadratic padding there is a second
way to locate the overlap: give an explicit sufficient height for each
family.  This also explains why the same quartic polynomial appears in
both constructions.

Fix $d,\eta>0$ and take $q=\eta\epsilon^2$ in
\eqref{c:rat-domain}.  For $|x|,|y|\le1/2$, put
\begin{equation}\label{c:rat-threshold-function}
 \Theta_x(y)=\epsilon\alpha xy
       +\epsilon^2\bigl(\alpha x^2c(y)-P(y)\bigr).
\end{equation}
We prove the following sufficient condition for the first family:
\begin{equation}\label{c:rat-threshold-inclusion}
 (y,t)\in\mathcal W,\quad -\tfrac12\le t\le\tfrac12,
 \quad t\ge\Theta_x(y)
 \quad\Longrightarrow\quad
 (y,t)=F_x(Y,T)\text{ for some }(Y,T)\in D_{\epsilon,q}.
\end{equation}
All bounds below are uniform in the displayed variables, with $d,\eta$
fixed.

On the rectangle
$\mathcal R=[-1,1]\times[-1,\tfrac12+d/2]$ the field is smooth and
has a finite Lipschitz constant $L_d$.  For small $\epsilon$, the map
\[
 (Y,T)\longmapsto (y,t)-\epsilon xW_d(Y,T)
\]
is a contraction of $\mathcal R$ into itself for
$|y|\le1/2$, $-1/2\le t\le1/2$.
Indeed the target rectangle has positive margins inside $\mathcal R$,
the field is bounded there, and $\epsilon L_d/2<1$.
The inverse label therefore exists and is unique in $\mathcal R$.
The inverse preserves the wedge directly.  Indeed,
\begin{align*}
 -b/2+a&=-(r/2+Y)\frac{2Y+2d+r}{8(r+d)^2},\\
 -b/2-a&=(r/2-Y)\frac{-2Y+2d+r}{8(r+d)^2}.
\end{align*}
The displayed coefficients are bounded on $\mathcal R$.  The target wall
functions $(1/2-t)/2\pm y$ are therefore the corresponding label wall
functions $r/2\pm Y$ multiplied by positive factors when $\epsilon$
is small.  Nonnegativity of the target wall functions implies
nonnegativity of the label wall functions, including equality on either
wall.  Thus this inversion argument does not require the surjectivity
proof in Section~\ref{c:rat-inverse}.

For fixed $x,y$, the inverse height $T$ increases strictly with $t$.
In fact if the target heights differ by $\Delta t>0$, and $p,p'$
are their inverse labels, then
\[
 |(p'-p)-(0,\Delta t)|\le\frac{\epsilon L_d}{2}|p'-p|
 \le\frac{\epsilon L_d/2}{1-\epsilon L_d/2}\,\Delta t.
\]
The last coefficient is less than one for small $\epsilon$.
At $t=\Theta_x(y)$, the inverse equations first give
$Y=y+O(\epsilon)$ and $T=O(\epsilon)$.
Using $b(y,0)=\alpha y$ improves the latter bound to
$T=O(\epsilon^2)$.  Consequently
\[
 Y=y+\epsilon xc(y)+O(\epsilon^2),\qquad
 b(Y,T)=\alpha y+\epsilon\alpha xc(y)+O(\epsilon^2),
\]
and substitution in the height equation gives
\[
 T=-\epsilon^2P(y)+O(\epsilon^3)
   =-\epsilon^2P(Y)+O(\epsilon^3).
\]
The rectangle inversion also applies at the threshold if
$(y,\Theta_x(y))$ lies outside the wedge.  Boundedness of the field
on $\mathcal R$ gives $Y=y+O(\epsilon)$ uniformly over the entire
target height interval, not only at the threshold.  Thus changing from
the threshold label to any of these labels changes $P(Y)$ by
$O(\epsilon)$.  Combining this bound with monotonicity gives, for all
$t\ge\Theta_x(y)$ under consideration,
\[
 T\ge-\epsilon^2P(Y)-C\epsilon^3
       \ge-\epsilon^2P(Y)-\eta\epsilon^2,
 \qquad T\ge-\tfrac14
\]
when $\epsilon$ is sufficiently small.
The inverse label is in the wedge, proving
\eqref{c:rat-threshold-inclusion}.

Applying the same argument with $x$ replaced by $-y$ gives a sufficient
condition $-t\ge\Theta_{-y}(x)$ for the second family.  But
\[
 \Theta_x(y)+\Theta_{-y}(x)
   =-\epsilon^2\alpha c_2(x^2-y^2)^2\le0.
\]
At least one of those two conditions holds for every point of $K$.
This proves coverage again.  The fixed positive $\eta$ absorbs the
uniform cubic inverse error; the area calculation in
Section~\ref{c:rat-area} then shows how small $\eta$ must be chosen
before fixing $\epsilon$.
